\begin.Abstract}
Multi-agent systems that support recursive spawning of descendant agents face a critical failure mode we term \emph{autonomous drift}: the gradual uncoupling of spawned agents from the authority and intent of their originators, ultimately resulting in behavior that is ungovernable and untraceable. This drift arises not from adversarial intent but from the architectural absence of a persistent, verifiable lineage bond between each agent and its ancestors. As recursive spawning propagates across generations, the chain of accountability frays until no operational link remains to the root human principal.

We propose a structural remedy embedded at the agent architecture level: \emph{immutable parent pointers} that form an unbroken, append-only chain of lineage back to the root human authority for every agent in the system. Each agent stores a cryptographically signed reference to its parent that cannot be modified, retroactively replaced, or severed without destroying the agent's operational identity. This invariance makes autonomous drift architecturally impossible, not merely probabilistically unlikely. Any attempt to spawn an agent outside the chain produces an identity that the system can immediately detect and reject, while any agent whose chain has been corrupted is identifiable by each downstream caller through constant-time chain verification.

We formalize the immutable parent pointer protocol, prove its adequacy for drift prevention under standard adversarial models, and specify its integration with the BX3 Framework's accountability primitives. Evaluation against synthetic recursive spawning workloads demonstrates that the protocol incurs negligible runtime overhead while eliminating all classes of drift-enabling lineage attacks. Keywords: recursive spawning, immutable parent pointers, autonomous drift, multi-agent systems, accountability.
\end.Abstract}
